\section{Neighborhood Complexity}
\label{apn:nc}

In this appendix, we give the full proof of \Cref{thm:nei-comp-NIP}.
We begin by introducing the relevant notions.
\subsection{Preliminaries}
\label{apn:prelims}



\paragraph{Asymptotic notation.}
Throughout this section, $\subpoly(n)$ denotes an unspecified function $f\from\N\to\mathbb R_{\ge 0}$ such that for every fixed $\eps>0$ we have $f(n)\le n^\eps$ for large enough $n$. We write $\subpoly_k(n)$ when $f$ may additionally depend on a fixed parameter $k$, and similarly we write $\subpoly_{\CC}(n)$ when $f$ may depend on a fixed class of graphs $\CC$. (So this notation is the same as writing $n^{o(1)}$, $n^{o_k(1)}$, or $n^{o_{\CC}(1)}$. However, we introduce the new notation as to make the variables and parameters explicit.) Similarly, $\polylog(n)$ denotes an unspecified function bounded by $p(\log n)$ for some polynomial $p$, and $\polylog_k(n)$ allows dependence of the polynomial on a fixed parameter~$k$.


\paragraph{Graphs.}
A graph $G$ consists of a set $V$ of vertices and a set $E\subset {V\choose 2}$ of edges, which are denoted $V(G)$ and $E(G)$, respectively.
For a graph $G$ with disjoint subsets $A,B \subseteq V(G)$,  let $G[A,B]$ denote the bipartite graph $(A,B,E')$ that is \emph{semi-induced} by $A$ and $B$, where $E'$ contains those edges of $G$ with one endpoint in $A$ and one endpoint in $B$.

\paragraph{Structures.}
% We give a more formal definition of transductions than in \Cref{sec:prelims}. 
A graph is viewed as a relational structure, equipped with a binary relation symbol $E$ denoting adjacency.
A \emph{unary expansion} of a relational structure $A$ is obtained from $A$ by 
adding several unary relations (also called unary predicates) to~$A$.
%  A (unary) \emph{expansion} of such a structure $S$ by a unary predicate $U$ is a structure $S'$ 
% with the same domain and relations as $S'$,
% and additional unary predicates.
% A (first-order) \emph{transduction} $T$ is specified 
% by a formula $\phi(x,y)$ in the signature of graphs (a binary adjacency relation symbol $E$) extended with additional unary predicates.
% Given a graph~$G$, its output $T(G)$
% is the set of all graphs $H$ that can be obtained by first arbitrarily expanding $G$ with unary predicates occurring in $\phi$, then defining a graph with the same vertex set as $G$ and edges $ab \in {V(G)\choose 2}$ such that $\phi(a,b)\lor\phi(a,b)$ holds,
% and then taking an arbitrary induced subgraph $H$ of the obtained graph.
% For a graph class $\CC$, we define $T(\CC)\coloneqq \bigcup \setof{T(G)}{G\in\CC}$. 





% \medskip
% In general, we will proceed in steps, where in step $i$ we maintain 
% a set $A_i\subset A$ and 
% a partition $\cal P_i$ of a subset $B_i$ of $B$.
% We will maintain that  
% \begin{align}\label{eq:inv-ineq}
% \sum_{P\in\cal P_i} |\inducedSS{P}{A_i}|\ge |A|^{1+\delta}/\polylog_i(|A|).  
% \end{align}

% For $i=0$, $\cal P_0=\set{B}$, while 
% for $i>0$, $\cal P_i$ is a refinement of $\cal P_{i-1}$ such that for $P\in \cal P_i$, 
% if $P'$ is the part of $\cal P_{i-1}$ containing $P$,
% then $P$ is carved out inside $P'$
% by asking for a unique sampled neighbor $x\in A_i$ in the positive or negative merge graph of $P'/A_i$.
% In particular, $P$ is contained in the corresponding neighborhood.
% In particular, it will follow that $\inducedSS{P}{A_i}$ has VC-dimension at most $d-i$.


% After $d$ steps of the construction, we will 
% conclude  that for $P\in \cal P_d$,  $P/A_d$ has VC-dimension $0$,
% so $|P/A_d|=1$.
% Therefore, by \eqref{eq:inv-ineq}, we get that $$|\cal P_d|\ge |A|^{1+\delta}/\polylog_d(|A|).$$  
% As each of the involved positive/negative merge graph is transducible from $G$
% we can transduce a structure $H$ equipped with $d$ binary relations of out-degree $1$, with one part equal to $\cal P_d$ and the other part equal to $A$, and every vertex in $\cal P_d$ is uniquely determined by its $d$ outneighbors. As $H$ is transducible from a monadically dependent graph class by a fixed transduction and the underlying bipartite graph $H'$ is $K_{d+1,d+1}$-free, it follows that  $H'$ belongs to a nowhere dense graph class. Using the known bound for nowhere dense classes, we obtain that for every fixed $\varepsilon>0$, assuming $|A|$ is large enough, we have:
% $$|\cal P_d|\le |A|^{1+\varepsilon}.$$
% Putting together the two inequalities we obtain $|A|^{\delta}\le |A|^{\varepsilon}\cdot \polylog_d(|A|)$, where $\varepsilon$ is any fixed constant. This is a contradiction, if $|A|$ large enough.

\subsection{A simple lemma on VC-dimension}
\lemVCind*
\begin{proof}[Proof of \Cref{lem:VC-induction}]
It is enough to prove the claim for the family $\cal F_+$ of all $v$-positive sets in $\cal F$; the argument for the family of $v$-negative sets is symmetric.
If $\cal F_+=\emptyset$, then $\VCdim(\cal F_+)=-\infty$ and there is nothing to prove.

Suppose $\cal F_+$ shatters a set $X\subseteq V$.
Since every set in $\cal F_+$ contains $v$, we must have $v\notin X$.
For each $Y\subseteq X$, choose $F_Y\in\cal F_+$ with $F_Y\cap X=Y$.
As $F_Y$ is $v$-positive, there is a set $F'_Y\in\cal F$ such that $F'_Y\triangle F_Y=\set{v}$.
Then $F'_Y\cap X=Y$, while exactly one of $F_Y$ and $F'_Y$ contains $v$.
Hence
\[
\{F_Y\cap (X\cup\set{v}),\ F'_Y\cap (X\cup\set{v})\}=\set{Y\cup\set{v},\ Y}.
\]
Since this holds for every $Y\subseteq X$, the family $\cal F$ shatters $X\cup\set{v}$.
Therefore,
\(
\VCdim(\cal F)\ge |X|+1,
\)
so $\VCdim(\cal F_+)\le \VCdim(\cal F)-1$.
This proves that $\cal F_+$ has VC-dimension strictly smaller than~$\cal F$.
\end{proof}





\subsection{Exhibiting many Hamming edges}
We prove the following generalization of \Cref{lem:star0}.

\begin{restatable}%[\inAppendix]
  {lemma}{lemstar}\label{lem:star}
  Let $G=(A,B,E)$ be a bipartite graph with $|A|=n>1$, and let $\cal P$ be a partition of $B$.
  Suppose
  \[
    m:=\sum_{P\in \cal P}|\inducedSS{P}{A}|> 2|\cal P|.
  \]
  Then there is a set $X\subseteq A$ such that
  \[
    \sum_{P\in \cal P}|E(\hamming(\inducedSS{P}{X}))|
    \ge \frac{m}{2\ln n+2}.
  \]
\end{restatable}
\begin{proof}
  Let $H(n):=1+1/2+\ldots+1/n$ be the $n$th harmonic number.
  Initially, set $X:=A$.
  As long as there is a vertex $v\in X$ such that 
  \[
    \sum_{P\in\cal P}\left(|\inducedSS{P}{X}|-|\inducedSS{P}{(X-\set v)}|\right)
    \le \frac{m}{2|X|\cdot H(n)},
  \]
  remove $v$ from $X$ and repeat. Otherwise, if there is no such $v$, terminate.
  

  \begin{claim}
    At the end of the process, the set $X$ is nonempty.
  \end{claim}
  \begin{proof}
    At any point in the process, consider the sum
    \[
      \mu(X):=\sum_{P\in \cal P}|\inducedSS{P}{X}|.
    \]
    Initially, when $X=A$, this sum is equal to $m$ by definition.
    In each step of the process, when a vertex $v$ is removed from $X$, the sum decreases by
    \[
      \sum_{P\in \cal P}\Big(|\inducedSS{P}{X}|-|\inducedSS{P}{(X-\set v)}|\Big)
      \le \frac{m}{2H(n)}\cdot \frac{1}{|X|}.
    \]
    Hence, if we remove all $n$ elements from $X$, the sum decreases by at most
    \[
      \left(\frac{m}{2H(n)}\cdot \frac{1}{n}\right) + \left(\frac{m}{2H(n)}\cdot \frac{1}{n-1}\right) + \ldots + \left(\frac{m}{2H(n)}\cdot \frac{1}{1}\right) = \frac{m}{2H(n)}\cdot \sum_{i=1}^n \frac 1 i
      = \frac{m}{2}.
    \]
    This means $\mu(\emptyset)\geq \mu(A) - m/2 = m/2$.
    On the other hand we have
    \[
      \mu(\emptyset)
      =\sum_{P\in\cal P}|\inducedSS{P}{\emptyset}|
      =\sum_{P\in\cal P}1
      =|\cal P|
      <\frac{m}{2};
    \]
    a contradiction which proves the claim.
  \end{proof}

  For every part $P\in\cal P$ and every vertex $v\in X$, the quantity
  \[
    |\inducedSS{P}{X}|-|\inducedSS{P}{(X-\set v)}|
  \]
  is exactly the number of edges of $\hamming(\inducedSS{P}{X})$ whose endpoints
  differ on~$v$. Therefore every edge of every graph
  $\hamming(\inducedSS{P}{X})$ is counted exactly once when we sum over
  $v\in X$, and hence
  \begin{multline*}
    \sum_{P\in\cal P}|E(\hamming(\inducedSS{P}{X}))|
    = \sum_{v\in X}
    \sum_{P\in\cal P}\left(|\inducedSS{P}{X}|-|\inducedSS{P}{(X-\set v)}|\right)\\
    > \sum_{v\in X}\frac{m}{2|X|\cdot H(n)}
    = |X|\cdot \frac{m}{2|X|\cdot H(n)}
    = \frac{m}{2H(n)}
    \ge \frac{m}{2\ln(n)+2}.
  \end{multline*}
\end{proof}


\subsection{Proof of \Cref{thm:nei-comp-NIP}}

 \Cref{thm:nei-comp-NIP} will follow directly from the next lemma.
\begin{restatable}{lemma}{lemnc}\label{lem:nc}
  Let $\BB$ be a monadically dependent class of bipartite graphs $G=(A,B,E)$
  with no twins in $B$.
  Then \begin{equation}\label{eq:bip-version}
  |B|\le |A|\cdot \subpoly_{\BB}(|A|).
\end{equation}
\end{restatable}
We first argue that \Cref{lem:nc} implies 
\cref{thm:nei-comp-NIP}.

\begin{proof}[Proof of \Cref{thm:nei-comp-NIP}]
Let $\BB$ be the class of all bipartite graphs $G[A,B]$
such that $G\in \CC$ and $A,B\subseteq V(G)$ are disjoint,
and no two vertices in $B$ have equal neighborhoods in $A$.
As $\CC$ transduces $\BB$, the class $\BB$ is monadically dependent.
For a graph $G\in \CC$ and set $A\subseteq V(G)$, let $B\subset V(G)-A$ 
be maximal such that no two vertices in $B$ have equal neighborhoods in $A$.
Then
\[
  \Big|\setof{N_G(v)\cap A}{v\in V(G)}\Big|\le  \Big|\setof{N_G(v)\cap A}{v\in A}\Big|+\Big|\setof{N_G(v)\cap A}{v\in B}\Big|   
  \le |A|+|B|,
\]
 Applying \Cref{lem:nc} to $G[A,B]\in\BB$ yields the theorem.
\end{proof}

We now proceed to the proof of \Cref{lem:nc}.
The following notion encapsulates the invariant used in the inductive proof.
\begin{definition}
  Let $G=(A,B,E)$ be a bipartite graph and $k\in\N$.
  A \emph{$k$-sparsification} consists of:
  \begin{itemize}
    \item sets $A_0\subseteq A$ and $B_0\subseteq B$, and
    \item functions $f_1,\ldots,f_k\from B_0\to A$,
  \end{itemize}
  such that
  the mapping $b\mapsto (N(b)\cap A_0,f_1(b),\ldots,f_k(b))$ is an injection from $B_0$ to $2^{A_0} \times A^k$.
For such a $k$-sparsification we define:
\begin{itemize}
  \item The \emph{associated partition} $\cal P$ as the partition of $B_0$ such that two vertices $b,b'$ are in the same part of $\cal P$ if and only if 
  $f_j(b)=f_j(b')$ for all $j=1,\ldots,k$. Note that vertices in the same part of $\cal P$ have pairwise distinct neighborhoods in $A_0$;
  \item The \emph{size} as $s=|B_0|$;
  \item The \emph{dimension} as $\max_{P\in\cal P} \VCdim(\inducedSS{P}{A_0})$; and
  \item The \emph{complexity} as the least $c\ge0$ such that there are first-order formulas $\phi_1(x,y),\ldots,\phi_k(x,y)$ of quantifier rank ${\le}c$, each using ${\le}c$ unary predicates, and there is an expansion $G'$ of $G$ with $c$ unary predicates,
  such that $\phi_j$ defines $f_j$ in $G'$ for $j=1,\ldots,k$, so that 
  \[
    f_j(b)=a\iff G'\models \phi_j(b,a)
    \quad
    \text{for all $b\in B_0$ and $a\in A$.}  
  \]
\end{itemize}
We say that a sparsification is \emph{terminal} if its size $s$ and
associated partition $\cal P$ satisfy $s\le 2|\cal P|$, and
\emph{nonterminal} otherwise.
\end{definition}

Note that a bipartite graph $G=(A,B,E)$ with no twins in $B$ has a trivial $0$-sparsification with $A_0=A$ and $B_0=B$. Its associated partition is $\cal P=\set{B}$, size $s=|B|$, and complexity is $0$.
We prove the following, for bipartite graphs $G$ from a monadically dependent class $\BB$:
\begin{enumerate}
  \item a nonterminal $k$-sparsification can be improved to a $(k+1)$-sparsification with strictly smaller dimension, by losing only a $\polylog(|A|)$ factor in the size, and increasing the complexity only by a constant (\Cref{lem:step});
  \item repeating this argument, we reach a terminal sparsification after at most $d$ steps, where $d$ upper bounds the VC-dimension of every $G$ in $\BB$ (\Cref{lem:iterate});
  \item  a terminal $k$-sparsification of $G\in \BB$ of bounded complexity has size $s\le |A|\cdot \subpoly_{\BB,k}(|A|)$ (\Cref{lem:terminal}).
\end{enumerate}
Combining these three points yields a terminal $k$-sparsification with $k\le d$ and size $s$ satisfying
$$\frac{|B|}{\polylog_d(|A|)}\le s\le |A|\cdot\subpoly_{\BB,d}(|A|).$$
This proves the inequality \eqref{eq:bip-version} in \Cref{lem:nc}. We now prove the necessary lemmas.


\begin{restatable}%[\inAppendix]
  {lemma}{lemtransduce}\label{lem:transduce}
  For every $k,c\in\N$ there is a transduction $T_{k,c}$ with the following property.
  Let $G=(A,B,E)$ be a bipartite graph with a $k$-sparsification 
   of complexity $c$, consisting of sets 
   $A_0\subseteq A,B_0\subseteq B$ and functions $f_1,\ldots,f_k\from B_0\to A$.
   Consider the bipartite graph $H=(A,B_0,E')$
with $N_H(b)=\setof{f_j(b)}{j\in[k]}$ for $b\in B_0$.
Then $H\in T_{k,c}(G)$.
\end{restatable}
\begin{proof}
Fix $k,c\in\N$.
Up to logical equivalence, there are only finitely many binary formulas of
quantifier rank at most $c$ over the signature consisting of the adjacency
relation and $c$ unary predicates.
Hence there are only finitely many $k$-tuples of such formulas; list them as
\[
  \bar\phi^1,\ldots,\bar\phi^m,
  \qquad
  \bar\phi^t=(\phi^t_1,\ldots,\phi^t_k).
\]

Let $T_{k,c}$ be the transduction that guesses:
\begin{itemize}
  \item $c$ unary predicates for the witness expansion from the definition of complexity,
  \item two unary predicates $L,R$ marking the sets $A$ and $B_0$, and
  \item unary predicates $S_1,\ldots,S_m$ used as global flags selecting one
  of the tuples $\bar\phi^1,\ldots,\bar\phi^m$.
\end{itemize}
It then applies the formula
\[
  \psi(x,y)\ :=\ L(x)\land R(y)\land
  \bigvee_{t=1}^m\left(\exists z\,S_t(z)\land \bigvee_{j=1}^k\phi_j^t(y,x)\right),
\]
where the inner disjunction is interpreted as false when $k=0$,
and finally takes the induced subgraph on $L\cup R$.

Now suppose $G$ comes with a $k$-sparsification of complexity $c$ as in the
statement.
By definition of complexity, after choosing a suitable expansion of $G$ by
$c$ unary predicates, there is some index $t\in[m]$ such that
$\phi_1^t,\ldots,\phi_k^t$ define the functions $f_1,\ldots,f_k$ on~$B_0$.
Interpret $L$ as $A$, $R$ as $B_0$, interpret $S_t$ as $V(G)$, and interpret
all other flags $S_{t'}$ as empty.
For this choice of unary predicates, the graph produced by $\psi$ has exactly
the edges $\{a,b\}$ with $a\in A$, $b\in B_0$, and $a=f_j(b)$ for some $j\in[k]$.
Therefore, the induced subgraph on $A\cup B_0$ is precisely the graph
$H=(A,B_0,E')$ with
\[
  N_H(b)=\setof{f_j(b)}{j\in[k]}
  \qquad\text{for every }b\in B_0.
\]
Hence, $H\in T_{k,c}(G)$.
\end{proof}

\begin{lemma}\label{lem:terminal}
  Let $\BB$ be a monadically dependent class of bipartite graphs and let $k,c\in\N$.
  Let $G=(A,B,E)\in\BB$ have a terminal $k$-sparsification of complexity at most $c$ and size $s$.
  Then
  \[
    s\le |A|\cdot \subpoly_{\BB,c,k}(|A|).
  \]
\end{lemma}
\begin{proof}
Let $A_0\subseteq A$, $B_0\subseteq B$, and $f_1,\ldots,f_k\from B_0\to A$
witness the given $k$-sparsification,
and let $\cal P$ be the associated partition. We prove that $|\cal P|\le |A|\cdot \subpoly_{\BB,c,k}(|A|)$, which will imply the claim as $s\le 2|\cal P|$ by terminality.
Choose a set $B^*\subseteq B_0$ containing exactly one vertex from each part of the associated partition
$\cal P$.

By \Cref{lem:transduce}, there is a transduction $T_{k,c}$ such that the
bipartite graph $H_0$ with parts $A$ and $B_0$ and
edges $\setof{\set{b,f_j(b)}}{b\in B_0,j\in[k]}$ 
belongs to $T_{k,c}(G)$.
Since transductions allow taking induced subgraphs, the graph
$H\coloneqq H_0[A\cup B^*]$ also belongs to $T_{k,c}(G)$.

Each vertex of $B^*$ has degree at most $k$ in $H$, so $H$ is
$K_{k+1,k+1}$-free.
Moreover, distinct vertices of $B^*$ lie in different parts of $\cal P$, hence
their tuples
\[
  (f_1(b),\ldots,f_k(b))
\]
are pairwise distinct.
Therefore, for every set $S\subseteq A$, there are at most $|S|^k$ vertices
$b\in B^*$ with $N_H(b)=S$, because each coordinate of the above tuple must
then belong to $S$.
Since every such $S$ has size at most $k$, there are in fact at most
$k^k$ such vertices (with $0^0=1$). Consequently, 
\[
  |\cal P|=|B^*|
  \le k^k\cdot \Big|\setof{N_H(b)}{b\in B^*}\Big|.
\]

Since $T_{k,c}(\BB)$ is again monadically dependent and $H\in T_{k,c}(\BB)$,
applying \Cref{thm:weaklySparseCase} with $t=k+1$ yields
\[
  \Big|\setof{N_H(b)}{b\in B^*}\Big|
  =\Big|\setof{N_H(b)\cap A}{b\in B^*}\Big|
  \le |A|\cdot \subpoly_{\BB,c,k}(|A|).
\]
The two inequalities,
together with $s\le 2|\cal P|$,
prove the
statement.
\end{proof}

The next lemma is the key ingredient of the proof of \Cref{thm:nei-comp-NIP}.
\begin{lemma}\label{lem:step}
  Suppose $G=(A,B,E)$ has a nonterminal $k$-sparsification of size $s$,
  dimension $d$, and complexity $c$.
  Then there is a $(k+1)$-sparsification of $G$ with:
  \begin{itemize}
    \item size at least $\Omega(\frac{s}{d\cdot \log^2(|A|)})$,
    \item dimension at most $d-1$,
    \item complexity at most $c+5$.
  \end{itemize}
\end{lemma}
\begin{proof}
Let $A_0\subseteq A$, $B_0\subseteq B$, and $f_1,\ldots,f_k\from B_0\to A$
witness the given $k$-sparsification.
Let $\cal P$ be its associated partition.
Since the sparsification is nonterminal, we have $s>2|\cal P|$.
For every part $P\in\cal P$, the functions $f_1,\ldots,f_k$ are constant on $P$,
so the injectivity condition in the definition of a sparsification implies that
the map $b\mapsto N_G(b)\cap A_0$ is injective on $P$.
Therefore,
\[
  |\inducedSS{P}{A_0}|=|P|
  \qquad\text{for every }P\in\cal P,
\]
and hence
\[
  \sum_{P\in\cal P}|\inducedSS{P}{A_0}|=\sum_{P\in\cal P}|P|=|B_0|=s.
\]
In particular, $d\ge 1$, since otherwise every set system $\inducedSS{P}{A_0}$
would have VC-dimension $0$ and hence size $1$, contradicting
$s>2|\cal P|$.
In particular, $|A_0|>1$, since otherwise every set system $\inducedSS{P}{A_0}$
would have size at most $2$, contradicting $s>2|\cal P|$.

Apply \cref{lem:star} to the bipartite graph $G[A_0,B_0]$ and the partition
$\cal P$ of $B_0$. Since
\[
  \sum_{P\in\cal P}|\inducedSS{P}{A_0}|=s>2|\cal P|,
\]
we obtain a set $A_1\subseteq A_0$ such that
\begin{equation}\label{eq:step-many-hamming-edges}
  \sum_{P\in\cal P}|E(\hamming(\inducedSS{P}{A_1}))|
  \ge \frac{s}{2\ln |A_0|+2}.
\end{equation}

Choose a set $B^*\subseteq B_0$ so that for each part $P\in\cal P$,
\[
  \inducedSS{(P\cap B^*)}{A_1}=\inducedSS{P}{A_1}
\]
and no two vertices in $P\cap B^*$ have equal neighborhoods in $A_1$.
Thus, the map $b\mapsto N_G(b)\cap A_1$ is a bijection from $P\cap B^*$ to
$\inducedSS{P}{A_1}$.

For each $P\in\cal P$, let $H_P$ be the Hamming graph of
$\inducedSS{P}{A_1}$, with vertex set identified with $P\cap B^*$ (identified using the previously constructed bijection between the two).
Let $H^*$ be the disjoint union of the graphs $H_P$, for $P\in\cal P$.
Then $V(H^*)=B^*$ and, by \eqref{eq:step-many-hamming-edges},
\[
  |E(H^*)|
  =\sum_{P\in\cal P}|E(H_P)|
  \ge \frac{s}{2\ln |A_0|+2}.
\]

For each $P\in\cal P$, let $M_P^+$ and $M_P^-$ denote the positive and
negative merge graphs of $\inducedSS{P}{A_1}$.
Let $M^+$ and $M^-$ be the disjoint unions of the graphs
$M_P^+$ and $M_P^-$, respectively, identified along the common part $A_1$.
As $A_1\subset A_0$,
% \By \Cref{obs:vc-hereditary},
 we have $\VCdim(\inducedSS{P}{A_1})\le \VCdim(\inducedSS{P}{A_0})\le d$ for every $P\in\cal P$.
Then \Cref{cor:hamming} implies that $H_P$ has at least $|E(H_P)|/d$
non-isolated vertices.
Summing over $P\in\cal P$, we infer that $H^*$ has at least
\[
  \frac{|E(H^*)|}{d}
  \ge \frac{s}{d(2\ln |A_0|+2)}
\]
non-isolated vertices.
Every such vertex is non-isolated in at least one
of $M^+$ and $M^-$, so for some $\sigma\in\set{+,-}$ the graph $M^\sigma$
has at least
\begin{equation}\label{eq:step-many-nonis}
  \frac{s}{2d(2\ln |A_0|+2)}
\end{equation}
non-isolated vertices in $B^*$.

We now choose a large subset on which each vertex has a unique
$M^\sigma$-neighbor.
We apply \cref{lem:dreier-sampling} to the bipartite graph obtained
from $M^\sigma$ by removing the isolated vertices. % from $B^*$.
We obtain sets $A_1'\subseteq A_1$ and $B_1\subseteq B^*$ such that every
vertex of $B_1$ has exactly one neighbor in $A_1'$ in $M^\sigma$, and
\[
  |B_1|
  \ge \Omega\left(\frac{s}{d\cdot \log^2|A|}\right),
\]
because $|A_1|\le |A_0|\le |A|$.

Define $f_{k+1}\from B_1\to A_1'$ by mapping each $b\in B_1$ to its unique
neighbor in $A_1'$ in the graph $M^\sigma$.
Let $\cal P_1$ be the family of all nonempty sets of the form
$f_{k+1}^{-1}(v)\cap P$, where $P\in\cal P$ and $v\in A_1'$.

We claim that $A_1$, $B_1$, and the functions
$f_1,\ldots,f_k,f_{k+1}$ form a $(k+1)$-sparsification of $G$.
Indeed, if $b,b'\in B_1$ lie in different parts of $\cal P$, then
$(f_1(b),\ldots,f_k(b))\neq (f_1(b'),\ldots,f_k(b'))$.
If they lie in the same part $P\in\cal P$, then 
$N_G(b)\cap A_1\neq N_G(b')\cap A_1$ whenever $b\neq b'$.
This is because we have chosen $B^* \supseteq B_1$ such that all vertices in $P \cap B^*$ have pairwise different neighborhoods on $A_1$ in $G$.
Hence, the map
\[
  b\mapsto (N_G(b)\cap A_1,f_1(b),\ldots,f_k(b),f_{k+1}(b))
\]
is injective on $B_1$.
Moreover, by construction, $\cal P_1$ is exactly the associated partition of
this new sparsification.

We now verify the claimed properties.
First, let $P'\in\cal P_1$.
Then $P'=f_{k+1}^{-1}(v)\cap P$ for some $P\in\cal P$ and $v\in A_1'$.
The set system $\inducedSS{P'}{A_1}$ is a subfamily of the neighborhood of $v$
in the $\sigma$-merge graph of $\inducedSS{P}{A_1}$, so
\Cref{lem:blue-vc} yields
\[
  \VCdim(\inducedSS{P'}{A_1})\le \VCdim(\inducedSS{P}{A_1}) - 1 \le \VCdim(\inducedSS{P}{A_0}) - 1 \le d-1,
\]
where the second inequality holds as $A_1\subset A_0$.
Thus, the new sparsification has dimension at most $d-1$.
Second, the size of the new sparsification is
\[
  s'\coloneqq |B_1|,
\]
and by the choice of $B_1$ we have
\[
  s' \ge \Omega\left(\frac{s}{d\cdot \log^2|A|}\right).
\]

Finally, let $G'$ be an expansion of $G$ witnessing that the original
$k$-sparsification has complexity $c$,
and let $\phi_1,\ldots,\phi_k$ be formulas defining
$f_1,\ldots,f_k$ in $G'$.
Expand $G'$ further by five unary predicates marking the sets
$A_1$, $A_1'$, $B_1$, $B^*$, and a set $S$ which is all of $V(G)$ if
$\sigma=+$ and empty if $\sigma=-$.
Using the formulas $\phi_1,\ldots,\phi_k$, the relation of belonging to the
same part of $\cal P$ is first-order definable with quantifier rank at most
$c+1$.
Consequently, the edge relation of $M^\sigma$ on $B^*\times A_1$ is also
first-order definable with quantifier rank at most $c+2$:
for $b\in B^*$ and $a\in A_1$, we ask whether there exists
$b'\in B^*$ in the same part of $\cal P$ such that
$N_G(b)\cap (A_1-\set{a})=N_G(b')\cap (A_1-\set{a})$ and
$b,b'$ differ on $a$ in the direction prescribed by $\sigma$.
Restricting this relation to $B_1\times A_1'$ gives a formula defining the
function $f_{k+1}$, of quantifier rank at most $c+2$.
Thus the new sparsification has complexity at most $c+5$.
(We only used quantifier rank $c + 2$, but $c+5$ many unary predicates.)
\end{proof}


\begin{lemma}\label{lem:iterate}
  Let $G=(A,B,E)$ be a bipartite graph with no twins in $B$, with
  $|A|>1$, and with $B\neq\emptyset$. 
  Suppose $\VCdim(\inducedSS{B}{A})\le d$ for some $d$.
  Then for some $i\le d$, the graph $G$ has a terminal $i$-sparsification 
  of complexity at most $5i$ 
  and size $s$ such that
  \[
    s\ge \frac{|B|}{\polylog_d(|A|)}.
  \]
\end{lemma}
\begin{proof}
  Set $n\coloneqq|A|$ and $m\coloneqq|B|$.
Let $\mathbf S_0$ be the trivial $0$-sparsification of $G$, given by
$A_0\coloneqq A$ and $B_0\coloneqq B$.
Since there are no twins in $B$, this is indeed a sparsification, with
size $s_0=m$, dimension at most $d$, and complexity $c_0=0$.

For $j=0,1,2,\ldots$, construct $\mathbf S_j$ recursively as follows.
Write $s_j$, $d_j$, and $c_j$ for the size, dimension, and
complexity of $\mathbf S_j$, and write $\cal P_j$ for its associated
partition.

If $\mathbf S_j$ is terminal, or if $j=d$, stop and set $i\coloneqq j$.
Otherwise, \Cref{lem:step} yields a $(j+1)$-sparsification
$\mathbf S_{j+1}$ of size $s_{j+1}$ and complexity $c_{j+1}$ such that
\[
  s_{j+1}
  \ge \Omega\left(\frac{s_j}{d_j\cdot \log^2 n}\right)
  \ge \Omega\left(\frac{s_j}{d\cdot \log^2 n}\right),
\]
while $c_{j+1}\le c_j+5$ and $d_{j+1}\le d_j-1$.

Let $0 \leq i \leq d$ be the index for which the above construction stops.
Repeated application of the above bounds
gives
\[
  s_i\ge \frac{m}{\polylog_d(n)}
  \qquad\text{and}\qquad
  c_i\le 5i\qquad\text{and}\qquad d_i\le d-i.
\]
We argue that $\mathbf S_i$ is terminal,
which will yield the conclusion for $s\coloneqq s_i$.

If $i<d$, then $\mathbf S_i$ is terminal by construction.
Suppose now that $i=d$.
Then $\mathbf S_d$ has dimension~$0$.
Let $P\in\cal P_d$.
By definition of the associated partition, the functions
$f^d_1,\ldots,f^d_d$ are constant on $P$.
Since $\mathbf S_d$ is a sparsification, the map
$b\mapsto N_G(b)\cap A_d$ is injective on~$P$. 
Hence, the vertices of $P$ have pairwise different neighborhoods on $A_d$, and we have $|\inducedSS{P}{A_d}| = |P|$.
Furthermore, as $d_i=0$, the set system $\inducedSS{P}{A_d}$ has VC-dimension $0$,
and hence size $1$.
Thus, also $|P|=1$ for every $P\in\cal P_d$, and therefore
\[
  s_d=|B_d|=|\cal P_d|\le 2|\cal P_d|.
\]
So $\mathbf S_d$ is terminal as well.
\end{proof}


We now prove \cref{lem:nc}, which we restate here for convenience.

\lemnc*

\begin{proof}
Fix now a graph $G=(A,B,E)\in\BB$,
and let $n\coloneqq |A|$ and $m\coloneqq |B|$.
If $|A|\le 1$ or $|B|=0$ then the statement is trivial, so suppose $|A|\ge 2$ and $|B|\ge 1$.
As $\BB$ is monadically dependent, the set system $\inducedSS{B}{A}$ has
VC-dimension bounded by a constant $d$ depending only on~$\BB$.
By \Cref{lem:iterate}, for some $i\le d$, the graph $G$ has a
terminal $i$-sparsification of size $s$ and complexity at most $5i$ such
that
\[
  s\ge \frac{m}{\polylog_d(n)}.
\]
By \Cref{lem:terminal}, as $i\le d$, we have
\[
  s
  \le n\cdot \subpoly_{\BB,d}(n).
\]
Combining the two inequalities, we obtain
\[
  m
  \le n\cdot \polylog_d(n)\cdot \subpoly_{\BB,d}(n)\le n\cdot \subpoly_{\BB}(n),
\]
 which proves  \eqref{eq:bip-version}.
\end{proof}
